\section{Weighted Chernoff Bound Derivation} \label{sec:appendix}


\begin{theorem}[(Weighted) Chernoff bound]
    Let $X_1, X_2, \dots, X_n$ be independent random variables taking values in $\sbr{0,1}$, let $w_1, w_2, \dots, w_n$ be non-negative weights, such that for every $i\in \sbr{n}$ we have that $w_i\in \sbr{0,1}$. Let $X=\sum_{i=1}^{n} w_i\cdot X_i$ and let $\mu = \E{X}$. Then for every $\delta_1 > 0$ and $\delta_2 \in (0,1)$ we have that:
    \[
        \Pr\sbr{X \geq (1+\delta_1)\cdot \mu} \leq \br{\frac{e^{\delta_1}}{(1+\delta_1)^{1+\delta_1}}}^\mu
        \qquad
        \Pr\sbr{X \leq (1-\delta_2)\cdot \mu} \leq \br{\frac{e^{-\delta_2}}{(1-\delta_2)^{1-\delta_2}}}^\mu.
    \]
\end{theorem}


\begin{proof}
We first prove the upper-tail bound for $\delta_1>0$. For any $\lambda>0$, Markov's inequality gives:
\[
\Pr\sbr{X \ge \br{1+\delta_1}\cdot\mu}
= \Pr\sbr{e^{\lambda X} \ge e^{\lambda\cdot\br{1+\delta_1}\cdot\mu}} \le \frac{\E{e^{\lambda X}}}{e^{\lambda\cdot\br{1+\delta_1}\cdot\mu}}.
\]
By independence, we have $
\E{e^{\lambda \cdot X}}
= \prod_{i=1}^n \E{e^{\lambda\cdot w_i\cdot X_i}}$.
Fix $i\in\sbr{n}$. Since $w_i\cdot X_i\in\sbr{0,1}$, by convexity of $y\mapsto e^{\lambda \cdot y}$ on $\sbr{0,1}$, for every $y\in\sbr{0,1}$,
$e^{\lambda \cdot y}
\le \br{1-y}\cdot e^0 + y\cdot e^{\lambda}
= 1 + \br{e^{\lambda}-1}\cdot y$.
Substituting $y=w_i\cdot X_i$, taking expectation, and using $1+t\le e^t$, we obtain:
\[
\E{e^{\lambda\cdot w_i\cdot X_i}}
\le 1 + \br{e^{\lambda}-1}\cdot w_i\cdot\E{X_i} \le \exp\br{\br{e^{\lambda}-1}\cdot w_i\cdot\E{X_i}}.
\]
Hence:
\[
\E{e^{\lambda X}}
\le \exp\br{\br{e^{\lambda}-1}\cdot\sum_{i=1}^n w_i\cdot\E{X_i}} = \exp\br{\br{e^{\lambda}-1}\cdot\mu}.
\]
Therefore,
$\Pr\sbr{X \ge \br{1+\delta_1}\cdot\mu}
\le \exp\br{\mu\cdot\br{e^{\lambda}-1-\lambda\cdot\br{1+\delta_1}}}$.
The function $f\br{\lambda}=e^{\lambda}-1-\lambda\cdot \br{1+\delta_1}$ is minimized at $\lambda=\ln\br{1+\delta_1}$, and so:
\[
\Pr\sbr{X \ge \br{1+\delta_1}\cdot\mu}
\le \exp\br{-\mu\cdot\br{\br{1+\delta_1}\cdot\ln\br{1+\delta_1}-\delta_1}} = \br{\frac{e^{\delta_1}}{\br{1+\delta_1}^{1+\delta_1}}}^{\mu}.
\]

Now let $\delta_2\in(0,1)$. For any $t>0$, again by Markov's inequality:
\[
\Pr\sbr{X \le \br{1-\delta_2}\cdot\mu}
= \Pr\sbr{e^{-tX} \ge e^{-t\cdot\br{1-\delta_2}\cdot\mu}} \le \frac{\E{e^{-t\cdot X}}}{e^{-t\cdot\br{1-\delta_2}\cdot\mu}}.
\]
Applying the same mgf estimate with $\lambda=-t<0$ yields $
\E{e^{-t\cdot X}}
\le \exp\br{\br{e^{-t}-1}\cdot\mu}$,
so:
\[
\Pr\sbr{X \le \br{1-\delta_2}\cdot\mu}
&\le \exp\br{\mu\cdot\br{e^{-t}-1+t\cdot\br{1-\delta_2}}}.
\]
The right-hand side is minimized at $t=-\ln\br{1-\delta_2}$, which gives:
\[
\Pr\sbr{X \le \br{1-\delta_2}\cdot\mu}
\le \exp\br{-\mu\cdot\br{\delta_2+\br{1-\delta_2}\cdot\ln\br{1-\delta_2}}} = \br{\frac{e^{-\delta_2}}{\br{1-\delta_2}^{1-\delta_2}}}^{\mu}.
\]
\end{proof}

\section{Preprocessing Details} \label{sec:appendix-preprocessing}

In this section, we provide the full preprocessing routine used for both the \coverage and \connectivity problems.

\subparagraph*{Scaling of the costs}
Firstly, we produce a polynomial number of instances, in such a way that for at least one of them we have a good lower bound on $\OPT$. For convenience, we allow costs to be rational numbers and scale them appropriately so that all of them lie in the interval $\sbr{0,1}$. To do so, by testing consecutive powers of $2$, we guess the smallest number $b$, such that the costliest activated interface among all vertices in the optimal solution has cost at most $2^b$. Let $C=\cl{\log\br{\max_{v\in V\br{G}, i\in \lambda\br{v}} \brc{c(i,v)}}}$. The number of the guesses is $C+1$, which is polynomial in the input size. If either some edge cannot be covered by interfaces of cost at most $2^b$, or there is no vertex $v$ and interface $i\in \lambda\br{v}$ such that $c\br{i,v}\in\sbr{1/2,1}$, then we discard such a guess. Otherwise, for each vertex $v\in V$, we temporarily discard from $\lambda\br{v}$ all interfaces with cost larger than $2^b$. Let $C_b=\max_{v\in V\br{G}, i\in \lambda\br{v}} \brc{c(i,v)}$ in this newly created instance. We normalize the costs of remaining interfaces by dividing them by $C_b$, so that the largest cost is exactly $1$. For the right guess of $b$, this gives $\OPT\geq C_b\geq 1/2$.

\subparagraph*{Cheap and expensive vertices}
Having established the scaling of costs, we now consider two types of vertices with respect to the new costs. For a vertex $v\in V$, if $\sum_{i\in \lambda\br{v}} c\br{i,v}\leq 1$, then we can safely assume that $v$ activates all interfaces in $\lambda\br{v}$, since this costs at most $1$ and is thus upper bounded by $2\cdot\OPT$. We call such vertices \emph{cheap}. The remaining vertices are \emph{expensive}. For every expensive vertex $v$, we can assume that for any output of the algorithm, $\COST\br{\lambda_A,v}\geq 1$, since enforcing this changes the cost by a factor of at most $2$ and does not affect asymptotic guarantees.

\subparagraph*{Finding the solution}
Since our algorithms are randomized, we also need to control the probability of failure. We ensure that, with high probability, for all values of $b$ we find a good solution. To do so, for each guess of $b$, we repeat the randomized procedure $K=\cl{\log_m\br{C}+1}$ times (which is polynomial). The algorithm is as follows.

\begin{algorithm}[H]
\caption{Preprocessing algorithm} \label{alg:preprocessing}
$\lambda_A(v) \gets \lambda(v)$ for every $v\in V$.

\For{$b\in \brc{0, \dots, C}$}
{
    For the remainder of this iteration, discard interfaces with cost larger than $2^b$ and divide all remaining costs by $\max_{v\in V\br{G}, i\in \lambda\br{v}} \brc{c(i,v)}$.

    \If{there is no solution for this new instance or there is no vertex $v$ and interface $i\in \lambda\br{v}$ such that $c\br{i, v}\in \sbr{1/2, 1}$}
    {
        \textbf{continue}
    }

    \For{$j\in \sbr{K}$}
    {
        $\lambda_A' \gets$ the output of the randomized procedure for the scaled-cost instance.

        \If{$\COST\br{\lambda_A'} < \COST\br{\lambda_A}$ and $\lambda_A'$ is a covering/connecting assignment}
        {
            $\lambda_A \gets \lambda_A'$.
        }
    }
}

\Return{$\{\lambda_A(v)\}_{v\in V}$.}
\end{algorithm}

Assuming that the failure probability of the randomized procedure is at most $O\br{1/m}$, for each fixed $b$ the probability of failure after $K$ repetitions is at most $O\br{1/m^K}=O\br{1/(mC)}$. Since there are $C$ guesses for $b$, a union bound implies that the probability that there exists some value of $b$ for which the algorithm fails is at most $O\br{1/m}$. In particular, this also holds for the correct guess of $b$.